\section{Review of Related Literature}

\subsection{StyleGAN Latent Editing and Inversion}

StyleGAN introduced a style-based generator in which an input latent is transformed into an intermediate $W$ space and injected into synthesis layers, yielding stronger disentanglement and image quality than earlier GAN designs \cite{karras2019stylebasedgeneratorarchitecturegenerative}. StyleGAN2 improved this generator by replacing AdaIN with weight modulation and demodulation and by regularizing the latent-to-image mapping, which supports higher fidelity and more stable inversion \cite{karras2020analyzingimprovingimagequality}. These architectures enable latent traversal, where an image can be edited by adding an attribute direction to a latent code. Methods such as InterFaceGAN, GANSpace, and SeFa show that pose, color, expression, and other attributes can often be controlled through approximately linear directions in StyleGAN spaces \cite{shen2020interpretinglatentspacegans,harkonen2020ganspacediscoveringinterpretablegan,shen2021closedformfactorizationlatentsemantics}.

The latent representation chosen for editing strongly affects the balance between detail and control. A compact $W$ code tends to stay closer to the generator distribution and supports smoother edits, while $W^{+}$ assigns a separate style vector to each synthesis layer and improves reconstruction at the risk of less predictable edits. Feature-space representations can retain still finer spatial information, but they require an editing mechanism that does not simply overwrite details during synthesis. This tradeoff is central to FeatureSAGE because few-shot generation needs both stable class-level edits and preservation of local cues such as flower petal structure or animal-face texture.

For real images, latent editing first requires GAN inversion. Encoder-based methods such as pSp, e4e, and ReStyle predict $W^{+}$ codes efficiently, while optimization-based methods can improve reconstruction at the cost of speed \cite{richardson2021encodingstylestyleganencoder,tov2021designingencoderstyleganimage,alaluf2021restyleresidualbasedstyleganencoder}. The central difficulty is the fidelity-editability tradeoff: high-capacity inversions in $W^{+}$ can preserve fine detail, but may move codes away from regions where semantic edits remain reliable \cite{liu2023delvingstyleganinversionimage}. StyleFeatureEditor addresses this limitation by predicting both a $W^{+}$ latent and an intermediate feature tensor, then training a Feature Editor to update feature maps during editing \cite{bobkov2024devildetailsstylefeatureeditordetailrich}. This motivates its integration into SAGE, whose baseline pSp inversion may lose local information that is useful for class-consistent attribute editing.

This distinction matters because a few-shot image generator is not judged only by whether an edited output looks realistic in isolation. It must also preserve the category cues that make the single reference useful. If inversion discards local structure before the edit is applied, the downstream attribute editing stage cannot recover those details reliably. Conversely, if the representation preserves details but becomes too entangled, edits may change category identity. FeatureSAGE is positioned between these needs by retaining SAGE's disentangled attribute dictionary while using SFE to carry additional feature-level evidence through the decoder.

\subsection{Few-Shot Attribute Group Editing}

Few-shot image generation seeks to synthesize diverse samples for a novel category from one or a few references. Directly adapting a powerful generator to a very small target set can overfit and reduce diversity, so many approaches transfer a pretrained generator and regularize adaptation through freezing, contrastive losses, or distance preservation \cite{mo2020freezediscriminatorsimplebaseline,ojha2021fewshotimagegenerationcrossdomain,zhao2023closerlookfewshotimage}. Editing-based methods avoid full retraining by reusing latent directions learned from seen categories.

AGE formulates a category as a combination of category-relevant structure and category-irrelevant attributes, learning a sparse dictionary of transferable directions and applying them to unseen-class embeddings \cite{AGE}. SAGE improves the stability of AGE by estimating a class embedding using all available few-shot references and by selecting attribute directions from similar seen categories, thereby reducing class-inconsistent edits \cite{ding2023stableattributegroupediting}. TAGE further studies trustworthy attribute group editing using codebooks and prompt-driven semantic modules \cite{zhang2024tagetrustworthyattributegroup}. These works show that latent attribute dictionaries can support few-shot synthesis without retraining the generator, but their effectiveness still depends on inversion quality and on preserving class identity under multi-attribute changes.

The present work therefore keeps the SAGE factorization and adaptive editing logic rather than introducing a new few-shot generator. Its modification is narrower: the inversion and decoding path is strengthened using SFE so that the latent edit is paired with an edited feature tensor. This scope preserves the original SAGE research line while testing whether the known reconstruction benefits of feature-level inversion translate to one-shot attribute group editing.

The related work also motivates the selected baselines. AGE is the direct attribute-group editing baseline because it demonstrates few-shot generation through category-relevant and category-irrelevant decomposition. SAGE is the stronger baseline because it stabilizes AGE using class embedding mapping and adaptive editing. Comparing FeatureSAGE against both methods isolates whether the proposed inversion and feature-editing path improves a mature editing-based pipeline rather than merely outperforming an older formulation.
